\documentclass[11pt,a4paper]{article}

\usepackage[a4paper,margin=2.5cm]{geometry}

\usepackage{fontspec}
\usepackage[spanish]{babel}


\title{Reporte de Práctica: Cómputo Paralelo Utilizando Hilos}
\author{Ingeniería en Inteligencia Artificial}
\date{}

\begin{document}

\maketitle
\section*{Resumen}

La práctica se centró en la investigación y aplicación del cómputo paralelo mediante el uso de la biblioteca de hilos en lenguaje C dentro de un entorno de sistema operativo Linux. Inicialmente, se analizó la teoría detrás de la gestión de procesos e hilos, comprendiendo su ciclo de vida, la forma en que comparten los recursos del sistema y la gestión del espacio de memoria. Posteriormente, se desarrollaron diferentes códigos de prueba para asimilar el comportamiento de las funciones básicas de la interfaz de programación. Se implementó un programa inicial que demostró la creación de un hilo simple que ejecutaba iteraciones concurrentes, mientras que el proceso padre esperaba su conclusión utilizando mecanismos de unión para sincronizar la finalización.

Luego de comprender la sincronización básica, se exploró la configuración de atributos para generar hilos desvinculados. Mediante un programa secundario, se instanciaron múltiples flujos de ejecución que operaron de forma totalmente autónoma, permitiendo que el proceso principal continuara y finalizara sin necesidad de esperar a sus procesos hijos. Asimismo, se investigó y comprobó rigurosamente la comunicación bidireccional entre procesos a través de apuntadores. Se codificó un bloque donde el proceso padre envió un arreglo de valores numéricos a un hilo; este último modificó dichos valores directamente en el espacio de memoria compartida, y el proceso padre recuperó e imprimió los datos actualizados exitosamente tras la conclusión del hilo.

Finalmente, como objetivo práctico integrador de la investigación, se desarrolló un programa enfocado en el procesamiento paralelo de estructuras de datos bidimensionales. El proceso principal generó una matriz numérica conformada por tres filas y nueve columnas. Para paralelizar la carga de trabajo, el proceso padre instanció tres hilos independientes, comunicándose con ellos mediante apuntadores para asignar a cada uno una fila específica de la matriz. La tarea designada para cada hilo consistió en procesar los datos recibidos, imprimir su segmento asignado y calcular el producto de los elementos de su respectiva fila. El proceso padre asumió el rol de coordinador, detectando el momento exacto en que los tres hilos terminaron sus respectivos cálculos para, posteriormente, recolectar los resultados alojados en la memoria y mostrarlos en la pantalla, evidenciando el éxito en la comunicación y distribución del procesamiento.
\clearpage
\section*{Conclusiones Personales}
\subsection{Said Carbot }
Durante la realización de esta práctica y la investigación previa sobre sistemas operativos y la gestión de procesos, pude comprender a profundidad la vital importancia del cómputo paralelo en la actualidad. Anteriormente, la concepción de la programación se limitaba a un enfoque estrictamente secuencial, donde cada instrucción debía esperar pacientemente a que la anterior concluyera su ciclo en el procesador. Sin embargo, al adentrarme en el desarrollo con la biblioteca de hilos en lenguaje C, experimenté un cambio de paradigma significativo. La investigación bibliográfica me permitió entender teóricamente qué es un hilo y cómo se diferencia fundamentalmente de un proceso completo, destacando que los hilos comparten el mismo espacio de memoria y descriptores de archivos, lo cual reduce drásticamente la sobrecarga de creación y el cambio de contexto por parte del sistema operativo subyacente. Esto se comprobó empíricamente al compilar y ejecutar los códigos de prueba, donde la instanciación de múltiples hilos desvinculados se realizó de manera fluida y con un consumo mínimo de recursos computacionales.

Al analizar la comunicación entre el proceso padre y los hilos mediante el uso de apuntadores, corroboré la naturaleza de la memoria compartida que detalla la literatura técnica. En uno de los ejercicios desarrollados, se enviaron argumentos a un hilo a través de un arreglo de enteros convertido a un apuntador genérico. Resultó fascinante observar cómo las modificaciones realizadas por el hilo sobre esos valores de memoria específicos se reflejaban de inmediato en el proceso padre una vez que el hilo terminaba su ejecución. Esta característica, aunque inmensamente poderosa para el paso de información, también me hizo reflexionar sobre los peligros inherentes de la concurrencia. Si múltiples hilos intentaran modificar la misma dirección de memoria simultáneamente sin los mecanismos de sincronización adecuados, se producirían condiciones de carrera. Aunque en esta práctica los ejercicios estaban estructurados lógicamente para evitar conflictos directos mediante la asignación de filas independientes, la investigación teórica dejó muy claro que el diseño cuidadoso y preventivo es indispensable para garantizar la integridad de los datos en sistemas paralelos reales.

El ejercicio principal, enfocado en el procesamiento de una matriz matemática de tres por nueve, consolidó todos los conceptos teóricos y prácticos abordados durante el periodo de estudio. El proceso padre asumió el rol de coordinador central, encargándose de inicializar los datos y distribuir la carga de trabajo de manera equitativa. Al crear tres hilos distintos y asignar a cada uno una fila específica para realizar la multiplicación continua de sus elementos, se logró una verdadera paralelización de tareas informáticas. Pude visualizar cómo la unidad central de procesamiento asignaba diferentes núcleos a cada hilo según la disponibilidad, permitiendo que las multiplicaciones se ejecutaran de forma simultánea.

En el contexto de la Ingeniería en Inteligencia Artificial, que es el marco principal de esta formación académica, los conocimientos adquiridos en esta práctica trascienden el mero ejercicio de programación. La inteligencia artificial moderna, especialmente el aprendizaje profundo y el entrenamiento de modelos masivos, depende fundamentalmente de la capacidad de paralelizar operaciones algebraicas a gran escala. Las multiplicaciones matriciales son el núcleo computacional de cualquier red neuronal. Entender desde el nivel más bajo cómo un sistema operativo gestiona estos recursos de procesamiento me ha dotado de una base teórica y práctica excepcionalmente sólida para optimizar algoritmos en el futuro. Me di cuenta de que el rendimiento real no solo depende de la complejidad teórica del código, sino de qué tan bien aprovecha la arquitectura física del procesador. La optimización del tiempo de ejecución demostrada en esta práctica es el primer peldaño para desarrollar sistemas inteligentes capaces de procesar información de manera instantánea. En conclusión, el desarrollo de la investigación y la implementación técnica de la práctica transformaron conceptos abstractos en habilidades tangibles que impactarán directamente en mi perfil profesional.
\clearpage
\end{document}